% Supplementary Material — Paper 2 (English). Built before manuscript.tex (cross-references).
\documentclass[11pt]{article}
\newcommand{\DocLang}{en}
\input{papers/shared/preamble.tex}
\input{papers/shared/authors.tex}
\usepackage{pdflscape}
\hypersetup{pdftitle={Supplementary Material: Lot release attributes and reactogenicity of VA-MENGOC-BC}}

\begin{document}
\SupplementaryNumbering
\SyntheticNotice
\begin{center}
{\Large\bfseries Supplementary Material\par}
\medskip
{\large Lot release attributes and reactogenicity of an outer membrane vesicle vaccine: linkage of
quality-control data from \PTNLots{} VA-MENGOC-BC lots with national adverse event reports\par}
\medskip
{\small \AuthorOne{} and \AuthorTwo\par}
\end{center}

\tableofcontents

\section{Control charts}\label{sec:S-charts}
\begin{table}[htbp]
\centering\small
\caption{Signals in the individuals chart (Nelson rules) and in the exponentially weighted moving
average chart ($\lambda=\PTEWMALambda$) of each attribute, in production order. Signals indicate
departures from a stable process, not non-conformity.}
\label{tab:S-charts}
\PTTableControlCharts
\end{table}

\section{National and export lots}\label{sec:S-equivalence}
\begin{table}[htbp]
\centering\small
\caption{Equivalence of national and export lots by two one-sided tests on the specification-window
scale (margin $\pm$\PTEqMarginPct{} of the window): difference in means with 90\% confidence
interval, and Hodges--Lehmann shift with 95\% confidence interval.}
\label{tab:S-equivalence}
\PTTableEquivalence
\end{table}

\begin{figure}[htbp]
\centering
\includegraphics{p2_equivalence}
\caption{Difference between national and export lots (90\% confidence interval) relative to the
equivalence margin (shaded).}
\label{fig:S-equivalence}
\end{figure}

\section{Linkage}\label{sec:S-linkage}
\begin{table}[htbp]
\centering\small
\caption{Reports involving VA-MENGOC-BC linked to a released lot, by year of vaccination.}
\label{tab:S-linkage}
\PTTableLinkage
\end{table}

\begin{table}[htbp]
\centering\small
\caption{Characteristics of reports linked and not linked to a released lot (assessment of linkage
bias).}
\label{tab:S-linkbias}
\PTTableLinkBias
\end{table}

\section{Sensitivity analyses}\label{sec:S-sensitivity}
\begin{landscape}
\begin{table}[htbp]
\centering\small
\caption{Odds ratios per SD (95\% CI) of the prespecified associations under alternative analyses:
primary GEE; restriction to exact lot links; exclusion of co-administered vaccines; retention of
temporally implausible links; Bayesian random-intercept logistic model (95\% credible interval);
lot-level quasi-binomial model of the proportion of reports with the reaction.}
\label{tab:S-sensitivity}
\PTTableSensitivity
\end{table}
\end{landscape}

\section{Exploratory associations}\label{sec:S-exploratory}
{\small
\PTTableGEEExploratory{Exploratory associations between every lot attribute (including multivariate
atypicality) and the primary and secondary reactions: logistic models with cluster-robust standard
errors by lot (independence working correlation), adjusted as in the primary analysis. $q$,
Benjamini--Hochberg adjusted $p$ within the exploratory family.\label{tab:S-exploratory}}}

\end{document}
